\documentstyle[%  


righttag%  


,11pt%  


,amssymb%  


,russian%               remove in Eng.  


,izvru%                 Set izven% in Eng.  


]{amsart}  


  


%       Math definitions  


 


%\newcommand{\ov}{\overline}  


\newcommand{\col}{\operatorname{col}}  


\newcommand{\vraisup}{\operatorname{vrai~sup}}  


  


\theoremstyle{plain}  


\theorembodyfont{\it}  


\newtheorem{thms}{\hskip\parindent Теорема}[section] 


\newtheorem{lems}{\hskip\parindent Лемма}[section] 


  


\theoremstyle{definition}  


\theorembodyfont{\rm}  


\newtheorem{notenonum}{\hskip\parindent Замечание}  


\renewcommand{\thenotenonum}{}  


\numberwithin{equation}{section} 


  


%\hoffset=-15mm%                  For Eng.  


  


\setcounter{page}{40}  


\god{1998}  


\nomer{2~(429)} %                Remove in Eng.  


%\nomer{Vol.\,42, No.\,2}%        Set in Eng.  


% \str{75--81}%                    Set in Eng.  


\udk{517.958:539.3}  


\author{\it В.А.\,Крысько, В.Н.\,Кузнецов, С.В.\,Полякова}  


% NB:   ^^ remove in Eng. 


\title{Операторный подход в геометрически\\ 


нелинейной задаче статической устойчивости\\ 


пластин и оболочек}  


  


\date{}  


% \pagestyle{myheadings}%         Set in Eng.  


  


\pagestyle{plain} %              Remove in Eng.  


\begin{document}  


  


% \thispagestyle{plain}%          Set in Eng.  


  


\maketitle  


 


\section*{Введение} 


В данной работе рассматривается задача определения допустимых 


значений проекций на оси функции напряжений $\sigma_x(x,y)$, 


$\sigma_y(x,y)$, $\tau_{xy}(x,y)$, действующих в серединной поверхности 


оболочки, 


при которой нелинейная система уравнений Маргерра-Власова-Муштари, 


а также система уравнений Тимошенко имеет 


единственное (устойчивое) решение. 


 


Одним из распространенных приемов решения задачи устойчивости 


пластин и оболочек на базе нелинейной модели является применение 


и реализация на ЭВМ численных методов \linebreak


(напр., [1]--[3]). В отдельных случаях (напр., [3]--[5]) задача 


решается путем выбора аппроксимирующих решений соответствующих 


уравнений. 


 


Авторы для решения данной задачи предлагают так называемый 


операторный подход, который сводит задачу устойчивости к 


задаче определения области параметров, при которых линейный 


оператор, связанный определенным образом с системой уравнений 


Маргерра-Власова-Муштари или с системой уравнений Тимошенко, 


является положительно определенным. 


 


\section{Основные положения операторного подхода\protect\\ 


в задаче устойчивости пластин и оболочек} 


 


Для удобства дальнейшего чтения кратко опишем основные положения 


операторного подхода, которые на примерах конкретных 


моделей будут подробно рассматриваться в следующих параграфах. 


 


Запишем систему уравнений Маргерра-Власова-Муштари или 


систему уравнений Тимошенко в операторной форме 


\begin{equation} 


A\ov u=\ov f, 


\end{equation} 


где $\ov u$, $\ov f$ --- элементы гильбертова пространства, выделяемого 


в пространстве $L_2(\Omega)\times L_2(\Omega)$ граничными условиями. В 


качестве 


параметров нелинейного оператора $A$ выступают значения функций 


$\sigma_x(x,y)$, $\sigma_y(x,y)$, $\tau_{xy}(x,y)$, которые в свою 


очередь зависят от 


величин нагрузок. Задача заключается в определении тех значений 


параметров, при которых уравнение (1.1) имеет единственное 


решение. 


 


Известно [6], что единственность решения уравнения (1.1) связана 


со свойством положительной монотонности оператора $A$. А 


именно, если оператор $A$ положительно монотонный, т.е. для 


любых $\ov u_1$, $\ov u_2$ выполняется неравенство 


$(A\ov u_1-A\ov u_2,\;\ov u_1-\ov u_2)\ge c\|\ov u_1-\ov u_2\|^2$, 


то уравнение (1.1) имеет не более одного решения. Тем самым 


задача устойчивости сводится к нахождению области допустимых 


значений параметров, при которых оператор $A$ уравнения 


(1.1) является положительно монотонным. Как будет показано 


ниже, эта задача в нашем случае сводится к задаче определения 


области допустимых значений параметров, при которых линейный 


оператор $A_1$, определенным образом связанный с оператором 


$A$, является положительно определенным в пространстве 


$\mu^2_0(\Omega)$, 


которое выделяется в пространстве $L_2(\Omega)$ однородными граничными 


условиями. 


 


\section{Применение операторного подхода\protect\\ 


в задаче устойчивости пластин и оболочек на базе\protect\\ 


модели Маргерра-Власова-Муштари} 


 


Запишем важную в теории оболочек модель Маргерра-Власова-Муштари 


[3] в операторной форме 


\begin{equation} 


A\ov u=(q,0), 


\end{equation} 


причем 


\begin{equation} 


A\ov u= 


\begin{cases} 


\frac{D}{h}\Delta^2w-L(w,F)-\Delta_kF,\\ 


\frac{1}{E}\Delta^2F+\frac{1}{2}L(w,w)+\Delta_kw,


\end{cases} 


\end{equation} 


для $\ov u=(W,F)\in H$, а 


гильбертово пространство $H$ выделяется в пространстве 


$L_2(\Omega)\times L_2(\Omega)$ граничными условиями вида 


\begin{equation} 


w|_\Gamma=\varphi_0,\quad 


\Bigl.\frac{\partial w}{\partial\eta}\Bigr|_\Gamma=\varphi_1,\quad 


F|_\Gamma=F_0,\quad 


\Bigl.\frac{\partial F}{\partial\eta}\Bigr|_\Gamma=F_1, 


\end{equation} 


где $\Gamma$ --- граница области $\Omega$. 


 


\begin{note} Геометрически нелинейную модель оболочки, 


определяемую операторным уравнением (2.1), часто в литературе 


(напр., [3]) называют моделью Кармана. Но основы геометрически 


нелинейной теории оболочек были заложены в [7], 


и свое завершение эта теория нашла в [8], [9]. 


Потому справедливости ради эту модель может быть следует называть 


моделью Маргерра-Власова-Муштари, именно такого названия модели 


авторы придерживаются в данной статье. 


\end{note} 


 


\begin{note} Вопросы существования решений в различных 


банаховых пространствах для систем, подобных нашей, хорошо 


изучены (напр., [10], [11]). Поэтому будем считать, что граничные 


условия и функция $q$ таковы, что модель Маргерра-Власова-Муштари 


имеет решение в пространстве $L_2(\Omega)\times L_2(\Omega)$. 


\end{note} 


 


Напомним хорошо известный факт [6]. 


 


\begin{lems} Если оператор $A$ вида {\em (2.2)} является положительно 


монотонным на множестве решений системы уравнений 


Маргерра-Власова-Муштари с граничными условиями {\em (2.3)}, то 


эта система уравнений имеет единственное решение. 


\end{lems} 


 


Далее, рассмотрим линейный оператор $A_1$, действующий в 


пространстве 


$$ 


H^2_0(\Omega)=\Bigl\{w\in L_2(\Omega);\;w|_\Gamma= 


\Bigl.\frac{\partial w}{\partial\eta}\Bigr|_\Gamma=0\Bigr\} 


$$ 


следующим образом: 


\begin{equation} 


A_1w=\frac{D}{h}\Delta^2w+\sigma_x 


\frac{\partial^2 w}{\partial x^2}+\sigma_y 


\frac{\partial^2 w}{\partial y^2}+2\tau_{xy} 


\frac{\partial^2 w}{\partial x\partial y}, 


\end{equation} 


где $\sigma_x$, $\sigma_y$, $\tau_{xy}$ --- проекции на оси функции 


напряжений, действующих 


в серединной поверхности оболочки в результате действия 


нагрузок. 


 


Справедлива 


 


\begin{thms} Пусть значения функций $\sigma_x$, $\sigma_y$, 


$\tau_{xy}$ таковы, 


что линейный оператор $A_1$ является положительно определенным. 


Тогда оператор $A$ вида {\em (2.2)} является положительно монотонным 


на множестве решений модели Маргерра-Власова-Муштари. 


\end{thms} 


 


{\bf Доказательство.} Нелинейное слагаемое оператора $A$ вида (2.2) 


обозначим через $B$. Тогда оператор $B$ определяется следующим 


образом: 


$$ 


B\ov u= 


\begin{cases} 


-L(w,F),\\ 


\frac{1}{2}L(w,w). 


\end{cases} 


$$ 


 


Пусть $\ov u_1$ и $\ov u_2$ --- два решения модели 


Маргерра-Власова-Муштари: 


$\ov u_1=(w_1,F_1)$, $\ov u_2=(w_2,F_2)$. Обозначим $v_1=w_1-w_2$, 


$v_2=F_1-F_2$. Тогда легко получается тождество 


\begin{equation} 


B\ov u_1-B\ov u_2= 


\begin{cases} 


-L(w_1,v_2)-L(v_1,F_2),\\ 


\frac{1}{2}L(v_1,v_1)+L(w_1,v_1). 


\end{cases} 


\end{equation} 


 


Для функций усилий $F_1$, $F_2$ имеем [3] 


\begin{equation} 


\frac{\partial^2 F_1}{\partial x^2}= 


\frac{\partial^2 F_2}{\partial x^2}=-\sigma_y;\quad 


\frac{\partial^2 F_1}{\partial y^2}= 


\frac{\partial^2 F_2}{\partial y^2}=-\sigma_x;\quad 


\frac{\partial^2 F_1}{\partial x\partial y}= 


\frac{\partial^2 F_2}{\partial x\partial y}=\tau_{xy}. 


\end{equation} 


Отсюда 


\begin{equation} 


\frac{\partial^2 v_2}{\partial x^2}= 


\frac{\partial^2 v_2}{\partial y^2}= 


\frac{\partial^2 v_2}{\partial x\partial y}=0 


\end{equation} 


и, следовательно, 


\begin{equation} 


L(w_1,w_2)=0. 


\end{equation} 


Рассмотрим оператор $T$, действующий в пространстве $H^2_0(\Omega)$, 


\begin{equation} 


Tv=T_x\frac{\partial^2 v}{\partial x^2}+2T_{xy} 


\frac{\partial^2 v}{\partial x\partial y}+T_y 


\frac{\partial^2 v}{\partial y^2}. 


\end{equation} 


Как показано в [12], при выполнении условий 


\begin{equation} 


\frac{\partial T_x}{\partial x}+ 


\frac{\partial T_{xy}}{\partial y}=0,\quad 


\frac{\partial T_{xy}}{\partial x}+ 


\frac{\partial T_y}{\partial y}=0 


\end{equation} 


имеет место равенство 


\begin{equation} 


(Tv_1,v_2)=(v_1,Tv_2). 


\end{equation} 


Преобразуем скалярное произведение (с учетом (2.5)) 


{\mathsurround=0pt


\begin{equation} 


(B\ov u_1\!-\!B\ov u_2,\ov u_1\!-\!\ov u_2)\!= \!


-(L(w_1,v_2),v_1)\!-\!(L(v_1,F_2),v_1)\!+\! 


\frac{1}{2}(L(v_1,v_1),v_2)\!+\!(L(w_1,v_1),v_2). 


\end{equation} 


}


Если на выражение $L(w,v)$ смотреть как на оператор вида (2.3), 


то легко проверить выполнение условий (2.10). Тогда в силу условий 


(2.1) и (2.7) получим 


\begin{equation} 


(L(w_1,v_1),v_2)+(L(v_1,v_1),v_2)=0, 


\end{equation} 


а в силу (2.6) 


\begin{equation} 


-L(v_1,F_2)=\sigma_x\frac{\partial^2 v_1}{\partial x^2}+ 


\sigma_y\frac{\partial^2 v_1}{\partial y^2}+2\tau_{xy} 


\frac{\partial^2 v_1}{\partial x\partial y}. 


\end{equation} 


Используя (2.8), (2.13), (2.14), запишем выражение (2.12) в виде 


\begin{equation} 


(B\ov u_1-B\ov u_2,\;\ov u_1-\ov u_2)=\biggl(\sigma_x 


\frac{\partial^2 v_1}{\partial x^2}+\sigma_y 


\frac{\partial^2 v_1}{\partial y^2}+2\tau_{xy} 


\frac{\partial^2 v_1}{\partial x\partial y},\;v_1\biggr). 


\end{equation} 


 


Учитывая (2.15), получим для другого скалярного произведения 


выражение 


\begin{multline} 


(A\ov u_1-A\ov u_2,\;\ov u_1-\ov u_2)= 


\biggl(\biggr[\frac{D}{h}\Delta^2v_1-\Delta_kv_2+\sigma_x 


\frac{\partial^2 v_1}{\partial x^2}+\\ 


+\sigma_y\frac{\partial^2 v_1}{\partial y^2}+2\tau_{xy} 


\frac{\partial^2 v_1}{\partial x\partial y}\biggr],\;v_1\biggr)+ 


\frac{1}{E}(\Delta^2v_2,v_2)+(\Delta_kv_1,v_2). 


\end{multline} 


Воспользуемся формулой Грина для операторов $\Delta$ и $\Delta_k$ 


(напр., [12]) 


и получим два соотношения: 


\begin{align} 


(\Delta^2v_1,v_1)&=\iint_\Omega\Delta^2v_1\cdot v_1d\Omega= 


\iint_\Omega\Delta v_1\cdot\Delta v_1d\Omega+\int_\Gamma 


\biggl(v_1\frac{\partial\Delta v_1}{\partial\eta}- 


\Delta v_1\frac{\partial\Delta v_1}{\partial\eta}\biggr)ds=\notag\\ 


&=\iint_\Omega\Delta v_1\cdot\Delta v_1d\Omega= 


(\Delta v_1,\Delta v_1),\\ 


(\Delta_k v_2,v_1)&=\iint_\Omega \Delta_k v_2\cdot v_1d\Omega= 


\iint_\Omega v_2\Delta_k v_1d\Omega+\int_\Gamma 


\biggl(v_1\frac{\partial v_2}{\partial\eta}-v_2 


\frac{\partial v_1}{\partial\eta}\biggr)ds=\notag \\ 


&=\iint_\Omega v_2\Delta_k v_1d\Omega=(v_2,\Delta_k v_1). 


\end{align} 


Учитывая (2.17), (2.18), для (2.16) будем иметь 


$$ 


(A\ov u_1-A\ov u_2,\;\ov u_1-\ov u_2)=(A_1v_1,v_1)+ 


\frac{1}{E}(\Delta^2v_2,v_2).


$$ 


Воспользуемся условием теоремы и тем фактом, что оператор $\Delta^2$ 


является положительно определенным в пространстве $H^2_0(\Omega)$ 


(напр., [12]). 


Тогда окончательно получим 


$$ 


(A\ov u_1-A\ov u_2,\;\ov u_1-\ov u_2)\ge c_1\|v_1\|^2 


+c_2\|v_2\|^2\ge c\|\ov u_1-\ov u_2\|^2, 


$$ 


что и завершит доказательство теоремы 2.1. 


\medskip 


 


Как следствие теоремы 2.1 и леммы 2.1 получается 


 


\begin{thms} Пусть допустимые значения функций $\sigma_x(x,y)$, 


$\sigma_y(x,y)$, $\tau_{xy}(x,y)$ таковы, что линейный оператор $A_1$, 


действующий 


в пространстве $H^2_0(\Omega)$, является положительно определенным. 


Тогда модель Маргерра-Власова-Муштари имеет 


единственное $($устойчивое$)$ решение. 


\end{thms} 


 


Отметим, что задача определения области допустимых значений 


функций $\sigma_x(x,y)$, $\sigma_y(x,y)$, $\tau_{xy}(x,y)$, 


при которых оператор $A_1$ 


является положительно определенным, рассматривалась в [12], 


где доказана 


 


\begin{thms} Пусть функции $\sigma_x(x,y)$, $\sigma_y(x,y)$, 


$\tau_{xy}(x,y)$ 


удовлетворяют условиям 


\begin{equation} 


\|\sigma_x\|<\frac{CD}{2h},\quad \|\sigma_y\|<\frac{CD}{2h},\quad 


\|\tau_{xy}\|<\frac{CD}{2h}, 


\end{equation} 


где $C$ --- наибольшая из констант 


в неравенстве Фридрихса: для $w\in H^2_0(\Omega)$ 


$$ 


(\Delta^2w, w)\ge C\iint_\Omega\biggl\{ 


\biggl(\frac{\partial w}{\partial x}\biggr)^2+ 


\biggl(\frac{\partial w}{\partial y} \biggr)^2\biggr\}dx\,dy. 


$$ 


 


Тогда оператор $A_1$ является положительно определенным. 


\end{thms} 


 


Отметим, что оценки (2.19) достаточно грубы. Задача более 


точного описания области положительной определенности оператора $A_1$ 


сводится к случаю, когда $\sigma_x(x,y)$, $\sigma_y(x,y)$, 


$\tau_{xy}(x,y)$ --- постоянные функции. 


В этой ситуации при наличии полной ортонормированной 


системы собственных функций оператора $A_1$ удается получить 


точное описание области положительной определенности 


этого оператора. С этой целью докажем следующую лемму. 


 


\begin{lems} Если линейный симметричный оператор $A$ имеет 


полную ортонормированную систему $\{U_n\}$ в качестве системы 


собственных функций с собственными значениями $\{\lambda_n\}$ и 


$\lambda_1<\lambda_2<\dotsb<\lambda_n<\dotsb$, то оператор $A$ 


является положительно 


определенным тогда и только тогда, когда $\lambda_1>0$. 


\end{lems} 


 


{\bf Доказательство.} Известно [12], что все собственные значения 


положительно определенного оператора $A$ являются положительными. 


 


Обратно (в предположении положительности собственных значений 


$\lambda_n$, $n=1,2,\dotsc$) имеем для 


$$ 


u=\sum^\infty_1(u,u_n)u_n


$$ 


соотношение 


$$ 


Au=\sum^\infty_1(Au,u_n)u_n=\sum^\infty_1(u,Au_n)u_n= 


\sum^\infty_1\lambda_n(u,u_n)u_n, 


$$ 


откуда 


$$ 


(Au,u)=\sum^\infty_1\lambda_n(u,u_n)^2\ge 


\lambda_1\sum^\infty_1(u,u_n)^2=\lambda_1\|u\|^2. 


\quad\square 


$$ 


 


Как следствие леммы 2.2 получается 


 


\begin{lems} Пусть оба линейных положительно определенных 


оператора $A_1$ и $A_2$ имеют одну и ту же полную ортонормированную 


систему функций $\{u_n\}$ в качестве системы собственных 


функций с собственными значениями $\{\lambda^i_n\}$, $i=1,2$, причем 


\begin{equation} 


0<\lambda^1_1/\lambda^2_1\le\dotsb\le\lambda^1_n/\lambda^2_n\le\dotsb 


\end{equation} 


и 


\begin{equation} 


0<\lambda^1_1-\lambda^2_1<\dotsb<\lambda^1_n-\lambda^2_n<\dotsb 


\end{equation} 


 


Тогда область положительной определенности оператора 


\begin{equation}  


A=A_1-\lambda A_2 


\end{equation}  


совпадает с множеством 


$$ 


\lambda<\lambda^1_1/\lambda^2_1. 


$$ 


\end{lems} 


 


{\bf Доказательство.} Легко видеть, что при условиях леммы 2.3 


система $\{u_n\}$ будет системой собственных функций оператора $A$ 


вида (2.22) с собственными значениями 


$\lambda_n=\lambda^1_n-\lambda\lambda^2_n$. В силу 


(2.20) и (2.21) для собственных значений $\lambda_n$ выполняются 


условия леммы 2.2, что и завершает доказательство леммы 2.3. 


 


\begin{notenonum} Результат леммы 2.2 позволяет описать 


область 


положительной определенности оператора 


$$ 


A=A_0-\lambda A_1-\mu A_2-\eta A_3, 


$$ 


где $A_0$, $A_1$, $A_2$, $A_3$ --- положительно определенные операторы, 


а $\lambda$, $\mu$, $\eta$ --- числовые параметры. 


\end{notenonum} 


 


 


\section{Применение операторного подхода в задаче устойчивости \protect\\ 


пластин и оболочек на базе модели Тимошенко} 


 


Модель Тимошенко является уточнением модели Маргерра-Власова-Муштари, 


записанной в перемещениях: наряду с функциями $u$, $v$, $w$ 


рассматриваются еще уточняющие функции $\psi_x$, $\psi_y$. 


Вывод уравнений Тимошенко приводится в [5] и там же рассматривается 


подстановка 


$$ 


\psi_x=-\biggl(\frac{\partial\varphi}{\partial x}+ 


\frac{\partial\theta}{\partial y}\biggr),\quad 


\psi_y=\frac{\partial\theta}{\partial x}- 


\frac{\partial\varphi}{\partial y}, 


$$ 


которая позволяет систему уравнений Тимошенко привести к более 


простому виду относительно функций $u$, $v$, $w$, $\theta$, $\varphi$. 


В этом случае систему уравнений Тимошенко удается записать в удобной 


для дальнейших исследований приведенной форме относительно 


функций $w$, $F$, $\theta$, $\varphi$, где $F$ --- функция усилий. 


Запишем эту систему 


уравнений в операторной форме 


$$ 


A\ov u=(q,0,0,0), 


$$ 


где оператор $A$ действует в пространстве $H$, выделенном из 


$L_2(\Omega)\times L_2(\Omega)\times L_2(\Omega)\times L_2$ 


граничными условиями вида 


\begin{equation} 


w|_\Gamma=w_1,\quad F|_\Gamma=F_1,\quad \Bigl. 


\frac{\partial F}{\partial\eta}\Bigr|_\Gamma=F_2,\quad 


\theta|_\Gamma=\theta_1,\quad \varphi|_\Gamma=\varphi_1, 


\end{equation} 


следующим образом: для $\ov u=(w,F,\theta,\varphi)\in H$ 


\begin{equation} 


A\ov u= 


\begin{cases} 


-\Delta w-\frac{2}{k^2_0(1-\mu)}L(w,F)- 


\frac{2}{k^2_0(1-\mu)}\Delta_k F- 


\frac{6k^2_0(1-\mu)}{h^2}(w-\varphi),\\ 


\frac{2}{k^2_0(1-\mu)}\Delta^2F+\frac{1}{k^2_0(1-\mu)} 


L(w,w)+\frac{2}{k^2_0(1-\mu)}\Delta_k w,\\ 


-\Delta\theta+\frac{12k^2_0}{h^2}\theta,\\ 


-\Delta\varphi-\frac{6k^2_0(1-\mu)}{h^2}(w-\varphi). 


\end{cases} 


\end{equation} 


 


Повторяя рассуждения, приведенные при доказательстве теоремы 2.1 


и теоремы 2.2 из \S\,2, получим два утверждения. 


 


\begin{thms} Пусть допустимые значения функций $\sigma_x(x,y)$, 


$\sigma_y(x,y)$, $\tau_{xy}(x,y)$ таковы, что линейный оператор 


$A_1$, действующий 


в пространстве $H^2_0(\Omega)=\{w\in L_2(\Omega);\;w|_\Gamma=0\}$, 


$$ 


A_1w=-\Delta w-\frac{6k^2_0(1-\mu)}{h^2}w+\frac{1}{k^2_0(1-\mu)} 


\biggl(\sigma_x\frac{\partial^2 w}{\partial x^2}+\sigma_y 


\frac{\partial^2 w}{\partial y^2}+2\tau_{xy} 


\frac{\partial^2 w}{\partial x\partial y}\biggr), 


$$ 


является положительно определенным. Тогда оператор $A$ вида 


{\em (3.2)} является положительно монотонным на множестве решений 


модели Тимошенко. 


\end{thms} 


 


\begin{thms} В условиях теоремы {\em 3.1} система уравнений 


Тимошенко 


при граничных условиях {\em (3.1)} имеет единственное $($устойчивое$)$ 


решение. 


\end{thms} 


 


В заключение отметим, что в данной работе рассматривались 


модели изотропных оболочек, что является несущественным 


с точки зрения применения операторного подхода. Кроме того, 


можно получить условия на функцию усилий $F$, при которых 


соответствующая модель имеет единственное решение. 


 


\begin{thebibliography}{99} 


 


\bibitem{Gri} 


Григолюк Э.И., Кабанов В.В. {\em Устойчивость оболочек}. -- М.: 


Наука, 1978. -- 360~с. 


\bibitem{Kr} 


Крысько В.А. {\em Нелинейная статика и динамика неоднородных 


оболочек}. -- Саратов: Изд-во Сарат. ун-та, 1976. -- 212~с. 


\bibitem{V1} 


Вольмир А.С. {\em Устойчивость деформируемых систем}. - М.: 


Наука, 1967. -- 984~с. 


\bibitem{V2} 


Вольмир А.С. {\em Гибкие пластины и оболочки}. -- М.: Гостехиздат, 


1949. -- 784~с. 


\bibitem{V3} 


Вольмир А.С. {\em Нелинейная динамика пластин и оболочек}. -- 


М.: Наука, 1972. -- 437~с. 


\bibitem{Lio} 


Лионс Ж.-Л. {\em Некоторые методы решения нелинейных краевых 


задач}. -- М.: Мир, 1972. -- 588~с. 


\bibitem{Mar} 


Marguerre K. {\em Theorie der gekr\"ummten Platte grosser


Form\"anderung} // Luftfahrforschung (Berlin). -- 1939. -- Bd.\,1. -- 


S.\,413--426. 


\bibitem{Vl} 


Власов В.3. {\em Основы дифференциальных уравнений общей теории 


оболочек} // ПММ. -- 1944. -- Т.\,8. -- \No\,2. -- С.\,109--140. 


\bibitem{Mu} 


Муштари Х.М., Галимов К.З. {\em Нелинейная теория упругих 


оболочек}. -- Казань: Таткнигоиздат, 1957. 


\bibitem{Vor} 


Ворович И.И. {\em 0 существовании решений в нелинейной 


теории оболочек} // Изв. АН СССР. Сер. матем. -- 1955. -- Т.\,19. 


-- \No\,4. -- C.\,173--186. 


\bibitem{Mor} 


Морозов Н.Ф. {\em 0 нелинейных колебаниях тонких пластин с 


учетом инерции вращения} // ДАН СССР. -- 1967. -- Т.\,176. -- \No\,3 


-- C.\,522--525. 


\bibitem{Mih} 


Михлин С.Г. {\em Вариационные методы в математической физике}. -- 


2-е изд. -- М.: Наука, 1970. -- 512~с. 


 


\end{thebibliography} 


 


\vskip 2cm 


 


\post{\it Саратовский технический}{\it Поступила}


\post{\it университет}{13.03.1995} 


 


\end{document} 


